%!TEX program = lualatex
% Arquivo: relatorio_pratica_8_oxigenio_dissolvido_winkler.tex

\documentclass[a4paper,12pt]{article}

% =========================================================
% Idioma e tipografia
% =========================================================
\usepackage[brazilian]{babel}
\usepackage{fontspec}
\setmainfont{Latin Modern Roman}
\usepackage{microtype}

% =========================================================
% Layout
% =========================================================
\usepackage[top=3cm,bottom=2.5cm,left=3cm,right=2.5cm]{geometry}
\usepackage{setspace}
\onehalfspacing
\usepackage{indentfirst}
\setlength{\parindent}{1.25cm}

% =========================================================
% Matemática, química e unidades
% =========================================================
\usepackage{amsmath,amssymb}
\usepackage[version=4]{mhchem}
\usepackage{siunitx}
\sisetup{
  locale = DE,
  output-decimal-marker = {,},
  per-mode = symbol,
  exponent-product = \cdot
}

% =========================================================
% Tabelas, figuras e listas
% =========================================================
\usepackage{booktabs}
\usepackage{array}
\usepackage{graphicx}
\usepackage{float}
\usepackage{enumitem}
\usepackage{xcolor}
\usepackage[most]{tcolorbox}

\setlist[itemize]{noitemsep,topsep=3pt}
\setlist[enumerate]{noitemsep,topsep=3pt}

% =========================================================
% Links
% =========================================================
\usepackage{hyperref}
\hypersetup{
  colorlinks=true,
  linkcolor=black,
  urlcolor=blue,
  citecolor=black,
  pdfauthor={Israel Carneiro; José Mauro Xavier Elsas; Miguel Gayani},
  pdftitle={Prática 8 -- Determinação de oxigênio dissolvido em águas pelo método de Winkler}
}

% =========================================================
% Dados experimentais
% =========================================================
\newcommand{\MassaMediaIodato}{0.0151}
\newcommand{\VolMedioPadronizacao}{15.7}
\newcommand{\ConcTiossulfato}{0.02697}
\newcommand{\VolAmostra}{50.0}
\newcommand{\VolMedioAmostra}{1.57}
\newcommand{\OxigenioDissolvido}{6.76}

\begin{document}

% =========================================================
% Capa
% =========================================================
\begin{titlepage}
\centering

{\large Universidade do Estado do Rio de Janeiro -- UERJ\par}
\vspace{0.3cm}
{\large Química VIII / Oceanografia\par}

\vspace{1.2cm}

{\LARGE\bfseries Prática 8\par}
\vspace{0.4cm}
{\Large\bfseries DETERMINAÇÃO DE OXIGÊNIO DISSOLVIDO\\
EM ÁGUAS PELO MÉTODO DE WINKLER\par}

\vspace{0.8cm}

\IfFileExists{uerj.jpg}{
  \includegraphics[width=0.28\linewidth]{uerj.jpg}
}{
  \begin{tcolorbox}[
    width=0.30\textwidth,
    colback=white,
    colframe=black!35,
    boxrule=0.5pt,
    arc=0mm
  ]
  \centering
  \vspace{1.0cm}
  \textit{Logotipo da UERJ}
  \vspace{1.0cm}
  \end{tcolorbox}
}

\vfill

\begin{flushright}
\textbf{Alunos:}\\
Israel Carneiro\\
José Mauro Xavier Elsas\\
Miguel Gayani\\[0.5cm]
\textbf{Professor:} Ítalo Vieira
\end{flushright}

\vspace{1cm}
Rio de Janeiro -- 2026

\end{titlepage}

% =========================================================
\section{Introdução}

O oxigênio dissolvido (OD) é um dos parâmetros mais importantes para a
avaliação da qualidade de ambientes aquáticos. Sua presença está associada
principalmente à transferência de oxigênio da atmosfera para a água e à
produção fotossintética. Por outro lado, a respiração dos organismos, a
oxidação de espécies químicas reduzidas e a decomposição da matéria orgânica
consomem oxigênio. Dessa forma, a concentração de OD resulta do balanço entre
processos de produção, transferência e consumo.

A solubilidade do oxigênio depende de fatores como temperatura, pressão,
salinidade e circulação da água. Em sistemas naturais, valores muito baixos
podem indicar consumo elevado de matéria orgânica e condições desfavoráveis
aos organismos aeróbios. Em águas que circulam por estruturas metálicas,
entretanto, o poder oxidante do oxigênio também pode favorecer processos de
corrosão.

Neste experimento, o oxigênio dissolvido foi determinado pelo método de
Winkler, uma análise volumétrica de oxidação-redução baseada em iodometria.
Inicialmente, adiciona-se sulfato manganoso à amostra em meio alcalino,
formando-se hidróxido de manganês(II):

\begin{equation}
\ce{MnSO4 + 2NaOH -> Mn(OH)2 + Na2SO4}.
\end{equation}

Na presença de oxigênio, o manganês é oxidado e forma um precipitado marrom.
A apostila representa essa etapa por:

\begin{equation}
\ce{2Mn(OH)2 + O2 -> 2MnO(OH)2}.
\end{equation}

Essa etapa é denominada \textit{fixação} do oxigênio, pois converte o
\ce{O2} dissolvido em uma espécie química que pode ser determinada
posteriormente. Após a acidificação e a adição de iodeto de potássio, ocorre
a liberação de iodo, cuja quantidade é estequiometricamente relacionada ao
oxigênio originalmente presente na amostra. As etapas podem ser representadas
por:

\begin{align}
\ce{MnO(OH)2 + 2H2SO4 &-> Mn(SO4)2 + 3H2O},\\
\ce{Mn(SO4)2 + 2KI &-> MnSO4 + K2SO4 + I2}.
\end{align}

O iodo produzido é titulado com tiossulfato de sódio:

\begin{equation}
\ce{I2 + 2S2O3^2- -> 2I- + S4O6^2-}.
\end{equation}

O amido é adicionado próximo ao ponto final e forma um complexo intensamente
azul com as espécies de iodo presentes. O desaparecimento dessa coloração
indica o término da titulação. Pela estequiometria global do método, um mol de
\ce{O2} corresponde a quatro mols de \ce{S2O3^2-}.

Antes da análise da amostra, a solução de tiossulfato foi padronizada com
iodato de potássio, empregado como padrão primário. Em meio ácido e na
presença de excesso de iodeto, o iodato produz triiodeto, que reage com o
tiossulfato segundo:

\begin{align}
\ce{IO3- + 8I- + 6H+ &-> 3I3- + 3H2O},\\
\ce{3I3- + 6S2O3^2- &-> 9I- + 3S4O6^2-}.
\end{align}

Consequentemente, um mol de \ce{KIO3} corresponde a seis mols de
\ce{Na2S2O3}, relação utilizada na determinação da concentração efetiva do
titulante.

% =========================================================
\section{Objetivos}

\begin{itemize}
  \item padronizar uma solução de tiossulfato de sódio empregando iodato de
        potássio como padrão primário;
  \item determinar a concentração de oxigênio dissolvido em uma amostra de
        água do mar pelo método de Winkler;
  \item relacionar a quantidade de tiossulfato consumida à quantidade de
        oxigênio presente na amostra;
  \item expressar o resultado de oxigênio dissolvido em
        \si{\milli\gram\per\litre}.
\end{itemize}

% =========================================================
\section{Material utilizado}

\begin{itemize}
  \item bureta de \qty{50.00}{\milli\litre};
  \item suporte universal e garra para bureta;
  \item balança analítica;
  \item Erlenmeyers;
  \item recipiente de aproximadamente \qty{500}{\milli\litre} com rolha;
  \item pipetas volumétricas e graduadas;
  \item proveta;
  \item béquer;
  \item frasco lavador com água destilada;
  \item papel-toalha.
\end{itemize}

% =========================================================
\section{Reagentes e soluções}

\begin{itemize}
  \item iodato de potássio, \ce{KIO3}, previamente dessecado;
  \item solução de tiossulfato de sódio, \ce{Na2S2O3}, de concentração
        nominal \qty{0.025}{\mole\per\litre};
  \item solução de iodeto de potássio, \ce{KI}, a 10\%;
  \item solução de ácido sulfúrico, \ce{H2SO4},
        \qty{0.5}{\mole\per\litre};
  \item solução de ácido sulfúrico a 50\%;
  \item solução de sulfato manganoso, \ce{MnSO4}, a 80\%;
  \item solução de hidróxido de sódio, \ce{NaOH}, a 30\%;
  \item suspensão de amido, empregada como indicador;
  \item água destilada;
  \item amostra de água do mar.
\end{itemize}

% =========================================================
\section{Procedimento experimental}

O procedimento descrito a seguir corresponde às orientações registradas no
quadro durante a aula. Quando há diferença entre essas anotações e o roteiro
impresso da apostila, foram mantidos os volumes e as concentrações indicados
em sala, por representarem a execução efetiva da prática.

\subsection{Padronização da solução de \ce{Na2S2O3}}

Pesou-se aproximadamente \qty{0.0150}{\gram} de \ce{KIO3}, que foi
transferido para um Erlenmeyer e dissolvido em cerca de
\qty{25}{\milli\litre} de água destilada. Adicionaram-se
\qty{5}{\milli\litre} de solução de \ce{KI} a 10\% e
\qty{10}{\milli\litre} de solução de \ce{H2SO4}
\qty{0.5}{\mole\per\litre}.

A solução formada, inicialmente castanha, foi titulada com \ce{Na2S2O3}
até adquirir coloração amarelo-clara. Próximo ao ponto final, adicionou-se o
indicador de amido e prosseguiu-se a titulação, gota a gota, até o
desaparecimento da coloração azul. As medições realizadas foram reunidas para
o cálculo da massa média de \ce{KIO3} e do volume médio de titulante.

\subsection{Preparo da amostra e fixação do oxigênio}

Um recipiente com rolha foi preenchido com aproximadamente
\qty{500}{\milli\litre} de água do mar. Adicionou-se
\qty{1}{\milli\litre} de solução de \ce{MnSO4} a 80\%, seguida de
\qty{1}{\milli\litre} de solução de \ce{NaOH} a 30\%. O recipiente foi
tampado e agitado, promovendo a formação do precipitado e a fixação do
oxigênio dissolvido.

Em seguida, foram acrescentados \qty{1}{\milli\litre} de solução de
\ce{KI} a 10\% e \qty{5}{\milli\litre} de solução de \ce{H2SO4} a 50\%.
O frasco foi novamente tampado e agitado cuidadosamente, até a dissolução do
precipitado e a liberação do iodo.

\subsection{Titulação da amostra}

Transferiram-se \qty{50.0}{\milli\litre} da amostra preparada para um
Erlenmeyer. A solução foi titulada com o \ce{Na2S2O3} padronizado até a
mudança de castanho-escuro para amarelo-claro. Adicionou-se então
\qty{1}{\milli\litre} de suspensão de amido, produzindo coloração azul, e
a titulação prosseguiu até a solução ficar incolor. Foram reunidos três
volumes de titulante, correspondentes às determinações realizadas pelos
integrantes do grupo, para o cálculo do volume médio.

% =========================================================
\section{Resultados e cálculos}

\subsection{Padronização da solução de \ce{Na2S2O3}}

Os dados empregados na padronização estão apresentados na
Tabela~\ref{tab:padronizacao}. A segunda leitura de volume, parcialmente
rasurada na ficha, foi reconstruída como \qty{13.5}{\milli\litre}, valor
compatível com a média de \qty{15.7}{\milli\litre} registrada no próprio
documento.

\begin{table}[H]
\centering
\caption{Dados empregados na padronização da solução de tiossulfato de
sódio.}
\label{tab:padronizacao}
\begin{tabular}{lccc}
\toprule
\textbf{Grandeza} & \textbf{I} & \textbf{II} & \textbf{Média}\\
\midrule
Massa de \ce{KIO3} (\si{\gram})
  & 0,0150 & 0,0152 & 0,0151\\
Volume de \ce{Na2S2O3} (\si{\milli\litre})
  & 17,9 & 13,5 & 15,7\\
\bottomrule
\end{tabular}
\end{table}

As médias foram calculadas por:

\begin{align}
\overline{m}_{\ce{KIO3}}
&=\frac{0.0150+0.0152}{2}
=\boxed{\qty{0.0151}{\gram}},\\[0.2cm]
\overline{V}_{\ce{Na2S2O3}}
&=\frac{17.9+13.5}{2}
=\boxed{\qty{15.7}{\milli\litre}}.
\end{align}

A massa molar empregada para o \ce{KIO3} foi
\qty{214.00}{\gram\per\mole}. Como um mol de iodato corresponde a seis
mols de tiossulfato:

\begin{equation}
n_{\ce{Na2S2O3}}=6n_{\ce{KIO3}}
=6\frac{m_{\ce{KIO3}}}{M_{\ce{KIO3}}}.
\end{equation}

Logo, a concentração padronizada foi:

\begin{align}
C_{\ce{Na2S2O3}}
&=
\frac{6\overline{m}_{\ce{KIO3}}}
{M_{\ce{KIO3}}\overline{V}_{\ce{Na2S2O3}}}\\
&=
\frac{6\left(\qty{0.0151}{\gram}\right)}
{\left(\qty{214.00}{\gram\per\mole}\right)
 \left(\qty{15.7e-3}{\litre}\right)}\\
&=\qty{0.02697}{\mole\per\litre}
\approx
\boxed{\qty{0.027}{\mole\per\litre}}.
\end{align}

\subsection{Volume médio gasto na titulação da amostra}

Os volumes de titulante consumidos nas três determinações da amostra estão
apresentados na Tabela~\ref{tab:amostra}.

\begin{table}[H]
\centering
\caption{Volumes de tiossulfato consumidos na titulação da amostra.}
\label{tab:amostra}
\begin{tabular}{lccc}
\toprule
\textbf{Dados da titulação} & \textbf{I} & \textbf{II} & \textbf{III}\\
\midrule
Volume da amostra (\si{\milli\litre})
  & 50,0 & 50,0 & 50,0\\
Volume de \ce{Na2S2O3} (\si{\milli\litre})
  & 1,6 & 1,9 & 1,2\\
\textbf{Volume médio de \ce{Na2S2O3} (\si{\milli\litre})}
  & \multicolumn{3}{c}{\textbf{1,57}}\\
\bottomrule
\end{tabular}
\end{table}

O volume médio foi:

\begin{align}
\overline{V}_{\ce{Na2S2O3}}
&=\frac{1.6+1.9+1.2}{3}\\
&=\qty{1.5667}{\milli\litre}
\approx\boxed{\qty{1.57}{\milli\litre}}.
\end{align}

\subsection{Concentração de oxigênio dissolvido}

Um mol de \ce{O2}, de massa molar \qty{32.00}{\gram\per\mole}, corresponde
a quatro mols de tiossulfato. Portanto, cada mol de \ce{Na2S2O3} consumido
equivale a:

\begin{equation}
\frac{\qty{32.00}{\gram\per\mole}}{4}
=\qty{8.00}{\gram\per\mole}
\end{equation}

de oxigênio. A concentração de OD pode ser calculada por:

\begin{equation}
OD\,(\si{\milli\gram\per\litre})=
\frac{
C_{\ce{Na2S2O3}}\,
V_{\ce{Na2S2O3}}\,
8.00\times1000
}{V_{\mathrm{amostra}}},
\end{equation}

com os volumes expressos em litros. Substituindo a concentração padronizada
e o volume médio:

\begin{align}
OD
&=
\frac{
\left(\qty{0.02697}{\mole\per\litre}\right)
\left(\qty{1.5667e-3}{\litre}\right)
\left(\qty{8.00}{\gram\per\mole}\right)
\left(\qty{1000}{\milli\gram\per\gram}\right)
}{\qty{50.0e-3}{\litre}}\\
&=\boxed{\qty{6.76}{\milli\gram\per\litre}}.
\end{align}

Para mostrar a variação entre as determinações, também foram calculados os
resultados individuais mantendo-se a mesma concentração padronizada.

\begin{table}[H]
\centering
\caption{Resultados individuais e médio de oxigênio dissolvido.}
\label{tab:od}
\begin{tabular}{lcc}
\toprule
\textbf{Determinação}
& \textbf{Volume de \ce{Na2S2O3} (\si{\milli\litre})}
& \textbf{OD (\si{\milli\gram\per\litre})}\\
\midrule
I   & 1,6  & 6,90\\
II  & 1,9  & 8,20\\
III & 1,2  & 5,18\\
\midrule
\textbf{Média} & \textbf{1,57} & \textbf{6,76}\\
\bottomrule
\end{tabular}
\end{table}

% =========================================================
\section{Discussão}

A padronização forneceu concentração de
\qty{0.02697}{\mole\per\litre} para o \ce{Na2S2O3}, valor próximo da
concentração nominal de \qty{0.025}{\mole\per\litre}. A diferença entre os
dois valores justifica a padronização prévia, pois soluções de tiossulfato
podem sofrer alterações durante o armazenamento. O cálculo baseado no
\ce{KIO3} permite utilizar a concentração efetivamente determinada, em vez
de assumir apenas o valor nominal.

O volume médio de tiossulfato consumido na análise da água do mar foi
\qty{1.57}{\milli\litre}, correspondendo a uma concentração média de
oxigênio dissolvido de \qty{6.76}{\milli\gram\per\litre}. Esse resultado
indica que a amostra continha quantidade mensurável de oxigênio disponível e
é compatível, em ordem de grandeza, com uma água aerada. A comparação mais
específica com a saturação de oxigênio exigiria informações não registradas,
como a temperatura, a salinidade e a pressão no momento da coleta.

Os resultados individuais variaram de
\qty{5.18}{\milli\gram\per\litre} a
\qty{8.20}{\milli\gram\per\litre}. Essa dispersão pode estar associada à
identificação visual do ponto final, ao pequeno volume de titulante consumido
e à leitura da bureta. Quando o volume gasto é reduzido, diferenças de poucas
décimas de mililitro representam uma fração relevante do resultado final.

O método também é sensível ao manuseio da amostra. A introdução de bolhas, o
contato excessivo com o ar antes da fixação ou a existência de espaço vazio
no recipiente podem alterar a quantidade de oxigênio dissolvido. A adição dos
reagentes e o fechamento do frasco devem, portanto, ser realizados com
cuidado. A dissolução incompleta do precipitado ou a demora entre a liberação
do iodo e a titulação também podem contribuir para erros.

O amido foi acrescentado apenas quando a solução se encontrava amarelo-clara.
Esse procedimento favorece a observação de uma viragem mais nítida e evita a
formação prematura de um complexo amido--iodo muito estável, que poderia
dificultar a identificação do ponto final.

O procedimento anotado no quadro apresentou algumas diferenças de volume e
concentração em relação ao roteiro impresso da apostila. Neste relatório foram
registradas as condições orientadas durante a aula, especialmente a utilização
de uma alíquota de \qty{50.0}{\milli\litre} e os volumes de reagentes indicados
no quadro. Essa escolha mantém a descrição coerente com a execução do
experimento e com os dados empregados nos cálculos.

% =========================================================
\section{Conclusão}

A solução de tiossulfato de sódio foi padronizada com iodato de potássio,
utilizando-se a relação estequiométrica de seis mols de \ce{S2O3^2-} para
cada mol de \ce{IO3-}. A partir da massa média de
\qty{0.0151}{\gram} de \ce{KIO3} e do volume médio de
\qty{15.7}{\milli\litre} de titulante, obteve-se:

\begin{equation*}
\boxed{
C_{\ce{Na2S2O3}}=\qty{0.02697}{\mole\per\litre}
\approx\qty{0.027}{\mole\per\litre}
}.
\end{equation*}

Na determinação pelo método de Winkler, as três titulações da amostra
consumiram \qty{1.6}{\milli\litre}, \qty{1.9}{\milli\litre} e
\qty{1.2}{\milli\litre} de solução padronizada, resultando em volume médio
de \qty{1.57}{\milli\litre}. Para uma alíquota de
\qty{50.0}{\milli\litre}, a concentração média calculada foi:

\begin{equation*}
\boxed{OD=\qty{6.76}{\milli\gram\per\litre}}.
\end{equation*}

O experimento demonstrou a aplicação de uma titulação iodométrica à análise
de água e permitiu relacionar quantitativamente o tiossulfato consumido ao
oxigênio originalmente dissolvido. A dispersão observada entre as
determinações evidencia a importância do controle da amostragem, da fixação
do oxigênio, da leitura da bureta e da identificação do ponto final.

% =========================================================
\section*{Referência}

LAGO, Dalva Cristina Baptista do; GOIS, Jefferson Santos de; LANGONE,
Marta Antunes Pereira; SENNA, Lilian Ferreira de.
\textit{Química VIII -- Oceanografia: aulas práticas}.
Rio de Janeiro: Universidade do Estado do Rio de Janeiro,
Instituto de Química, Departamento de Química Analítica, 2017.

\end{document}
